\chapter{Nuisance Samplers}
\label{ch:nuisance}

\newthought{A nuisance sampler is an inner step}, not a main sampler. While the main sampler
(Chapters~\ref{ch:basic-samplers}--\ref{ch:mcmc-3}) chooses the physics points, a nuisance sampler
runs \emph{around each of those points} to optimise a nuisance parameter. \Jarvis\ provides one:
\sampler{Profile1D}.

\section{Where it fits}
\label{sec:nui-where}

A nuisance sampler is configured inside \yamlkey{Sampling.Nuisance} of the \emph{main} sampler's
block---it augments the main scan rather than replacing it. For every point the main sampler
proposes, the nuisance step is run before the point's score is finalised.

\section{Profile1D}
\label{sec:s-profile1d}

\paragraph{How it works} \sampler{Profile1D} profiles a \emph{single} nuisance parameter over a
one-dimensional range using \emph{golden-section search}---an efficient bracketing method that
repeatedly narrows the interval toward the optimum without needing derivatives. At each main point
it searches the nuisance parameter for the value that minimises or maximises the nuisance
likelihood (\yamlkey{TargetMode}), within \yamlkey{MaxAttempt} evaluations, optionally subject to
\yamlkey{PassCondition} checks. The profiled value then feeds the main point's likelihood.

\paragraph{Configuration} Under \yamlkey{Sampling.Nuisance}:

\begin{description}[style=nextline, leftmargin=1.2em]
  \item[\yamlkey{Method}] must be \sampler{Profile1D}.
  \item[\yamlkey{Variables}] the nuisance parameter (in practice exactly one), with the usual
        \yamlkey{name}/\yamlkey{description}/\yamlkey{distribution} (typically \yamlkey{Flat} or
        \yamlkey{Log} with \yamlkey{min}/\yamlkey{max}).
  \item[\yamlkey{LogLikelihood}] the nuisance objective, as \yamlkey{name}/\yamlkey{expression}
        terms.
  \item[\yamlkey{TargetMode}] \code{min} or \code{max}---whether to minimise or maximise the
        nuisance objective.
  \item[\yamlkey{MaxAttempt}] the search budget (golden-section iterations).
  \item[\yamlkey{PassCondition}] a list of boolean expressions a point must satisfy (use \code{[]}
        for none).
\end{description}

\begin{yamlwin}
Sampling:
  Method: "Bridson"
  Radius: 0.25
  MaxAttempt: 30
  Variables: [ ... ]            # the physics parameters
  LogLikelihood:
    - {name: L_total, expression: "-0.5 * ((z - 100.0) / 10.0)**2"}
  Nuisance:
    Method: Profile1D
    Variables:
      - name: nu
        description: nuisance scale
        distribution: {type: Flat, parameters: {min: 0.8, max: 1.2}}
    LogLikelihood:
      - {name: L_nu, expression: "-0.5 * ((nu - 1.0) / 0.05)**2"}
    TargetMode: max
    MaxAttempt: 12
    PassCondition: []
\end{yamlwin}

\begin{jnote}
\sampler{Profile1D} profiles \emph{one} parameter per main point. To profile several nuisance
parameters, the recommended pattern is to keep the parameter of interest in the main scan and treat
each nuisance dimension explicitly; multi-dimensional nuisance profiling is not a single built-in
step in the current version.
\end{jnote}

\section{Summary}
\label{sec:nui-summary}

A nuisance sampler optimises nuisance parameters at each physics point instead of sampling them
jointly. \sampler{Profile1D} does this in one dimension by golden-section search, set up under
\yamlkey{Sampling.Nuisance} with a \yamlkey{TargetMode} and a search budget.
